\documentclass[11pt,a4paper]{article}
\usepackage{amsmath}
\usepackage{amssymb}
\usepackage{geometry}
\usepackage{hyperref}
\usepackage{listings}
\usepackage{xcolor}
\usepackage{graphicx}
\usepackage{float}

\geometry{margin=1in}

\lstset{
  language=Python,
  basicstyle=\ttfamily\small,
  breaklines=true,
  commentstyle=\color{gray},
  keywordstyle=\color{blue},
  stringstyle=\color{red},
  tabsize=2,
  frame=tb,
  backgroundcolor=\color{white!95!black}
}

\title{Complete Explanation of Architecture Diagrams\\Variables, Equations, and Data Flow}
\author{Detailed Guide with All Equations}
\date{}

\begin{document}

\maketitle
\tableofcontents
\newpage

\section{Notation and Abbreviations}

\subsection{Dimension Variables (Always Used)}

\begin{table}[H]
\centering
\begin{tabular}{|c|c|l|}
\hline
\textbf{Symbol} & \textbf{Value} & \textbf{Definition} \\
\hline
$B$ & 16 & Batch size (number of samples processed at once) \\
\hline
$T$ & 50 & Sequence length (time steps in load history) \\
\hline
$D$ & 64 & Hidden/embedding dimension (internal representation size) \\
\hline
$N$ & 8 & Number of input parameters (soil/pile characteristics) \\
\hline
$n$ & 20 & Number of yield surfaces (stiffness degradation stages) \\
\hline
\end{tabular}
\end{table}

\textbf{What these mean:}

\begin{itemize}
  \item \textbf{B = 16}: We process 16 samples simultaneously (a batch)
  \item \textbf{T = 50}: Each load history has 50 time steps
  \item \textbf{D = 64}: We represent everything internally as 64-dimensional vectors
  \item \textbf{N = 8}: Soil parameters include 8 different measurements
  \item \textbf{n = 20}: Foundation can degrade through 20 different stages
\end{itemize}

\section{Tensor Notation}

When we write a shape like $(B, T, D)$, it means:
\begin{itemize}
  \item \textbf{First dimension}: $B = 16$ items (batch size)
  \item \textbf{Second dimension}: $T = 50$ items (sequence length)
  \item \textbf{Third dimension}: $D = 64$ items (embedding dimension)
  \item \textbf{Total elements}: $16 \times 50 \times 64 = 51,200$ numbers
\end{itemize}

\section{DIAGRAM 1: SimpleTransformer Architecture}

This is the \textbf{complete end-to-end model}. Let me explain every box:

\subsection{INPUT STAGE}

\subsubsection{1. Load Sequence (Green box, top-left)}

\begin{equation}
\text{load\_sequence} \in \mathbb{R}^{B \times T \times 1} = \mathbb{R}^{16 \times 50 \times 1}
\end{equation}

\textbf{What is it?}
\begin{itemize}
  \item Time series of wind and wave loads acting on the OWT
  \item 16 different scenarios (batch size)
  \item 50 time steps for each scenario
  \item 1 scalar value per time step (the load magnitude)
\end{itemize}

\textbf{Example values:}
\begin{lstlisting}
load_sequence[0, :, 0] = [0.5, 0.7, 0.6, ..., 0.8]  # Sample 0
load_sequence[1, :, 0] = [0.3, 0.4, 0.5, ..., 0.6]  # Sample 1
...
# Each row is a different load history
\end{lstlisting}

\textbf{Physical interpretation:}
\begin{itemize}
  \item First time step: initial load value
  \item Middle time steps: how load changes during typhoon
  \item Last time step: final load value
\end{itemize}

\subsubsection{2. Soil Parameters (Green box, top-right)}

\begin{equation}
\text{soil\_params} \in \mathbb{R}^{B \times N} = \mathbb{R}^{16 \times 8}
\end{equation}

\textbf{What is it?}
\begin{itemize}
  \item Constant properties of the soil and pile (don't change over time)
  \item 16 different foundation scenarios
  \item 8 different measurements per scenario
\end{itemize}

\textbf{The 8 parameters typically include:}
\begin{itemize}
  \item $G_{\max}$: Maximum shear modulus of soil
  \item $s_u$: Undrained shear strength
  \item $EI$: Bending stiffness of pile
  \item $\phi$: Friction angle
  \item $d_p$: Pile diameter
  \item $h_p$: Pile length
  \item $\gamma_d$: Damping ratio
  \item $I_D$: Relative density
\end{itemize}

\textbf{Example values:}
\begin{lstlisting}
soil_params[0] = [8.5, 45.2, 1200, 32.5, 0.5, 30, 0.05, 0.65]
soil_params[1] = [7.2, 38.1, 1100, 30.1, 0.5, 30, 0.06, 0.55]
# Each row is a different soil profile
\end{lstlisting}

\subsection{PROCESSING STAGE 1: Load Processing (Pink box)}

\subsubsection{3. Embed Loads + Positional Encoding (Pink box, left)}

\textbf{Input:} $\text{load\_sequence} \in \mathbb{R}^{16 \times 50 \times 1}$

\textbf{Operation:}

\begin{equation}
\text{embedded\_loads} = W_e \cdot \text{load\_sequence} + \text{pos\_enc}
\end{equation}

Where:
\begin{itemize}
  \item $W_e$: Weight matrix for embedding (learnable) $(1 \times D) = (1 \times 64)$
  \item $\text{pos\_enc}$: Positional encoding (pre-computed, not learnable)
  \item Formula: $\text{pos\_enc}_{t,d} = \sin(t/10000^{2d/D})$ for even $d$, $\cos(...)$ for odd $d$
\end{itemize}

\textbf{What this does:}
\begin{itemize}
  \item Transforms 1D load values into 64D embeddings
  \item Adds positional information so transformer knows time order
  \item Example: load at $t=0$ gets different encoding than $t=49$
\end{itemize}

\textbf{Positional Encoding Formula:}

\begin{equation}
\text{PE}(t, 2i) = \sin\left(\frac{t}{10000^{2i/D}}\right), \quad \text{PE}(t, 2i+1) = \cos\left(\frac{t}{10000^{2i/D}}\right)
\end{equation}

Where:
\begin{itemize}
  \item $t$: time step (0 to 49)
  \item $i$: dimension index (0 to 31)
  \item $D$: total dimension (64)
\end{itemize}

\textbf{Why sine/cosine?}
\begin{itemize}
  \item Provides different pattern for each position
  \item Relative positions are always the same
  \item Allows model to learn ordinal relationships
\end{itemize}

\textbf{Output:} $\text{embedded\_loads} \in \mathbb{R}^{16 \times 50 \times 64}$

\textbf{Computation:}
\begin{lstlisting}
# Pseudo-code
for batch b in 0:16
  for time t in 0:50
    embedded_loads[b, t, :] = embed_layer(load_sequence[b, t, 0])
    embedded_loads[b, t, :] += positional_encoding[t, :]
\end{lstlisting}

\subsubsection{4. Output Shape}

\begin{equation}
\text{shape} = (B, T, D) = (16, 50, 64)
\end{equation}

This is shown in the diagram as the green box after "Embed Loads + Positional Encoding".

\subsection{PROCESSING STAGE 2: Foundation Stiffness (Right path)}

\subsubsection{5. Slot-Attention Module (Purple box)}

\textbf{Input:} $\text{soil\_params} \in \mathbb{R}^{16 \times 8}$

\textbf{This module has 3 internal steps:}

\paragraph{Step 1: Input Encoding}

\begin{equation}
\text{encoded\_params} = \tanh(W_{enc} \cdot \text{soil\_params} + b_{enc})
\end{equation}

Where:
\begin{itemize}
  \item $W_{enc}$: Learnable weight matrix $(8 \times 64)$
  \item $b_{enc}$: Learnable bias $(64)$
  \item $\tanh$: Activation function (squashes to $[-1, 1]$)
\end{itemize}

\textbf{Output:} $\text{encoded\_params} \in \mathbb{R}^{16 \times 64}$

\paragraph{Step 2: Cross-Attention (Slots attend to inputs)}

The module has 21 \textbf{learnable slots}:

\begin{equation}
\text{slots} \in \mathbb{R}^{21 \times 64}
\end{equation}

Cross-attention computes:

\begin{equation}
\text{Attention} = \text{softmax}\left(\frac{Q \cdot K^T}{\sqrt{D}}\right) \cdot V
\end{equation}

Where:
\begin{itemize}
  \item $Q$ (Query): slots $(21 \times 64)$ - "What do I need to know?"
  \item $K$ (Key): encoded params $(16 \times 64)$ - "What information exists?"
  \item $V$ (Value): encoded params $(16 \times 64)$ - "What are the values?"
  \item $\sqrt{D} = \sqrt{64} = 8$ - Scaling factor
\end{itemize}

\textbf{Computation for one sample:}

\begin{equation}
\text{scores}_{ij} = \frac{\sum_d Q_i[d] \cdot K_j[d]}{\sqrt{64}}
\end{equation}

\begin{equation}
\text{attention\_weights}_{ij} = \frac{\exp(\text{scores}_{ij})}{\sum_k \exp(\text{scores}_{ik})}
\end{equation}

\begin{equation}
\text{attended\_output}_i = \sum_j \text{attention\_weights}_{ij} \cdot V_j
\end{equation}

\textbf{Output:} $\text{attended\_slots} \in \mathbb{R}^{16 \times 21 \times 64}$

(16 batches, 21 slots, 64 dimensions, now each slot has seen the input parameters)

\paragraph{Step 3: Self-Attention (Slots talk to each other)}

Now slots communicate with each other:

\begin{equation}
\text{self\_attention} = \text{softmax}\left(\frac{Q \cdot K^T}{\sqrt{D}}\right) \cdot V
\end{equation}

Where $Q$, $K$, $V$ are all the attended slots.

\textbf{Purpose:}
\begin{itemize}
  \item Slot for $H_0$ learns from context of other slots
  \item Slots for $H_n$ don't duplicate each other
  \item Each slot specializes in one distinct role
\end{itemize}

\textbf{Output:} $\text{refined\_slots} \in \mathbb{R}^{16 \times 21 \times 64}$

\subsubsection{6. Stiffness Outputs (Green boxes after Slot-Attention)}

\textbf{From refined slots, we decode:}

\paragraph{From Slot 1: Initial Stiffness Matrix}

\begin{equation}
H_0 = \text{MLP}_{H0}(\text{refined\_slots}[:, 0, :])
\end{equation}

\textbf{Where:}
\begin{itemize}
  \item MLP: Multi-layer perceptron (3 layers)
  \item Input: refined\_slots for position 0 (16 × 64)
  \item Output: initial stiffness matrices (16 × 3 × 3)
\end{itemize}

\textbf{The MLP internally:}

\begin{equation}
h_1 = \text{ReLU}(W_1 \cdot \text{slot} + b_1)
\end{equation}

\begin{equation}
h_2 = \text{ReLU}(W_2 \cdot h_1 + b_2)
\end{equation}

\begin{equation}
H_0 = W_3 \cdot h_2 + b_3
\end{equation}

\textbf{Output:} $H_0 \in \mathbb{R}^{16 \times 3 \times 3}$

This is a 3×3 symmetric matrix for each of the 16 samples:

\begin{equation}
H_0[b] = \begin{bmatrix}
h_{11}^{(b)} & h_{12}^{(b)} & h_{13}^{(b)} \\
h_{12}^{(b)} & h_{22}^{(b)} & h_{23}^{(b)} \\
h_{13}^{(b)} & h_{23}^{(b)} & h_{33}^{(b)}
\end{bmatrix}
\end{equation}

\paragraph{From Slots 2-21: Stiffness Drops and Activation Gates}

\begin{equation}
(H_n, w_n) = \text{MLP}_{drops}(\text{refined\_slots}[:, n, :]) \quad \text{for } n = 1, \ldots, 20
\end{equation}

Each MLP outputs 10 values:
\begin{itemize}
  \item First 9 values: flattened 3×3 matrix $H_n$
  \item Last 1 value: activation gate (passed through sigmoid)
\end{itemize}

\begin{equation}
\text{output}[0:9] = H_n \text{ (3×3 matrix)}
\end{equation}

\begin{equation}
w_n = \sigma(\text{output}[9]) = \frac{1}{1 + e^{-\text{output}[9]}}
\end{equation}

Where $\sigma$ is the sigmoid function.

\textbf{Outputs:}
\begin{itemize}
  \item $H_{\text{drops}} \in \mathbb{R}^{16 \times 20 \times 3 \times 3}$
  \item $w_{\text{gates}} \in \mathbb{R}^{16 \times 20}$
\end{itemize}

\subsection{PROCESSING STAGE 3: Stiffness Evolution (Pink box)}

\subsubsection{7. Stiffness Evolution Module}

\textbf{Input:} $H_0, H_{\text{drops}}, w_{\text{gates}}$

\textbf{Formula:}

\begin{equation}
H_{\text{evolved}} = H_0 - \sum_{n=1}^{20} w_n \cdot H_n
\end{equation}

\textbf{Detailed computation for one sample $b$:}

\begin{equation}
H_{\text{evolved}}[b] = H_0[b] - \sum_{n=1}^{20} w_n[b] \cdot H_n[b]
\end{equation}

Where:
\begin{itemize}
  \item $w_n[b] \in [0, 1]$: Scalar weight for sample $b$, drop $n$
  \item $H_n[b] \in \mathbb{R}^{3 \times 3}$: Matrix for sample $b$, drop $n$
  \item Multiplication: scalar times matrix (element-wise scaling)
\end{itemize}

\textbf{Step-by-step for one sample:}

\begin{equation}
\text{degradation}[b] = w_1[b] \cdot H_1[b] + w_2[b] \cdot H_2[b] + \cdots + w_{20}[b] \cdot H_{20}[b]
\end{equation}

\begin{equation}
H_{\text{evolved}}[b] = H_0[b] - \text{degradation}[b]
\end{equation}

\textbf{Physical meaning:}
\begin{itemize}
  \item Initially ($t=0$): all $w_n = 0$, so $H_{\text{evolved}} = H_0$
  \item As loading increases: $w_n$ increases (sigmoid activation)
  \item More drops activated: $H_{\text{evolved}}$ decreases (stiffness degrades)
  \item Finally ($t=T$): $w_n \to 1$, $H_{\text{evolved}} \approx H_0 - \sum H_n$
\end{itemize}

\textbf{Output:} $H_{\text{evolved}} \in \mathbb{R}^{16 \times 3 \times 3}$

This represents the foundation stiffness after all degradation from the entire load history.

\subsection{PROCESSING STAGE 4: Concatenation (Pink box)}

\subsubsection{8. Concatenate Stiffness + Loads}

We need to combine:
\begin{itemize}
  \item Embedded loads: $(16, 50, 64)$
  \item Evolved stiffness: $(16, 3, 3)$
\end{itemize}

\textbf{Process:}

\begin{equation}
H_{\text{token}} = \text{embed}(H_{\text{evolved}})
\end{equation}

Convert the 3×3 matrix into a 64D token:

\begin{equation}
H_{\text{flattened}} = \text{vec}(H_{\text{evolved}}) \in \mathbb{R}^{16 \times 9}
\end{equation}

\begin{equation}
H_{\text{token}} = W_{h} \cdot H_{\text{flattened}} + b_h \in \mathbb{R}^{16 \times 64}
\end{equation}

\textbf{Concatenation:}

\begin{equation}
\text{concat\_input} = [\text{H\_token}; \text{embedded\_loads}] \in \mathbb{R}^{16 \times 51 \times 64}
\end{equation}

\textbf{Composition:}
\begin{itemize}
  \item Position 0: Stiffness token (16, 64)
  \item Positions 1-50: Load tokens (16, 50, 64)
  \item Total: 51 tokens per sample, 64D each
\end{itemize}

\textbf{Output:} $(B, T+1, D) = (16, 51, 64)$

\subsection{PROCESSING STAGE 5: Transformer Encoder}

\subsubsection{9. Transformer Encoder (Multi-Head Self-Attention)}

\textbf{Input:} $(16, 51, 64)$

\textbf{Operation:}

\begin{equation}
\text{encoder\_output} = \text{TransformerEncoder}(\text{concat\_input})
\end{equation}

\textbf{Internal computation (one layer):}

\begin{equation}
\text{attention} = \text{MultiHeadAttention}(Q, K, V)
\end{equation}

\textbf{With multiple heads:}

\begin{equation}
\text{head}_i = \text{softmax}\left(\frac{Q_i K_i^T}{\sqrt{D/h}}\right) V_i
\end{equation}

\begin{equation}
\text{MultiHeadAttention} = \text{Concat}(\text{head}_1, \ldots, \text{head}_4) \cdot W_o
\end{equation}

Where:
\begin{itemize}
  \item $h = 4$: number of heads
  \item $W_o$: output projection matrix
\end{itemize}

\textbf{Purpose:}
\begin{itemize}
  \item Each position "looks at" all other positions
  \item Learns which positions are related
  \item Example: Which loads matter given current stiffness?
\end{itemize}

\textbf{Then feed-forward:}

\begin{equation}
\text{FFN}(x) = \text{ReLU}(W_1 x + b_1) W_2 + b_2
\end{equation}

\textbf{Output:} $\text{encoder\_output} \in \mathbb{R}^{16 \times 51 \times 64}$

The 51 tokens now contain globally processed information.

\subsection{PROCESSING STAGE 6: Transformer Decoder}

\subsubsection{10. Transformer Decoder (Cross-Attention)}

\textbf{Input 1 (Memory):} $\text{encoder\_output} \in \mathbb{R}^{16 \times 51 \times 64}$

\textbf{Input 2 (Target):} $\text{encoder\_output}$ (same, for simplicity)

\textbf{Decoder Layer:}

\begin{equation}
\text{cross\_attn} = \text{MultiHeadAttention}(Q_{\text{tgt}}, K_{\text{mem}}, V_{\text{mem}})
\end{equation}

Where:
\begin{itemize}
  \item $Q_{\text{tgt}}$: queries from target (what do we want to generate?)
  \item $K_{\text{mem}}, V_{\text{mem}}$: keys and values from encoder memory
\end{itemize}

\textbf{With masking (causal attention):}

\begin{equation}
\text{scores}_{ij} = \begin{cases}
\frac{Q_i K_j^T}{\sqrt{D}} & \text{if } j \leq i \\
-\infty & \text{if } j > i
\end{cases}
\end{equation}

This ensures position $i$ only sees positions $j \leq i$ (not future).

\textbf{Output:} $\text{decoder\_output} \in \mathbb{R}^{16 \times 51 \times 64}$

\subsection{PROCESSING STAGE 7: Frequency Prediction (Pink box)}

\subsubsection{11. Frequency Head (MLP)}

\textbf{Input:} First token of decoder output

\begin{equation}
\text{freq\_token} = \text{decoder\_output}[:, 0, :] \in \mathbb{R}^{16 \times 64}
\end{equation}

(The stiffness token, which now has full context)

\textbf{MLP computation:}

\begin{equation}
h = \text{ReLU}(W_1 \cdot \text{freq\_token} + b_1)
\end{equation}

\begin{equation}
\text{freq\_pred} = W_2 \cdot h + b_2
\end{equation}

Where:
\begin{itemize}
  \item $W_1 \in \mathbb{R}^{64 \times 64}$: First weight matrix
  \item $W_2 \in \mathbb{R}^{64 \times 1}$: Output weight matrix (single value)
\end{itemize}

\textbf{Output:} $\text{freq\_pred} \in \mathbb{R}^{16}$

This is 16 scalar values, one predicted frequency per sample.

\section{DIAGRAM 2: Slot-Attention Flow Diagram}

This expands the purple "Slot-Attention Module" box from Diagram 1.

\subsection{Input Parameters}

\begin{equation}
\text{Input Parameters} \in \mathbb{R}^{16 \times 8}
\end{equation}

Example for one sample:
\begin{lstlisting}
[G_max, s_u, EI, phi, d_p, h_p, gamma, I_D]
[8.5, 45.2, 1200, 32.5, 0.5, 30, 0.05, 0.65]
\end{lstlisting}

\subsection{Input Encoder}

\textbf{Formula:}

\begin{equation}
\text{encoded} = \text{ReLU}(W_{\text{enc}} \cdot \text{input} + b_{\text{enc}})
\end{equation}

Where:
\begin{itemize}
  \item $W_{\text{enc}} \in \mathbb{R}^{8 \times 64}$: transforms from 8D to 64D
  \item $b_{\text{enc}} \in \mathbb{R}^{64}$: bias term
\end{itemize}

\textbf{ReLU activation:}

\begin{equation}
\text{ReLU}(x) = \max(0, x)
\end{equation}

\textbf{Output:} $\text{Encoded Input} \in \mathbb{R}^{16 \times 64}$

\subsection{Cross-Attention}

\textbf{Explanation:}

Each of 21 slots "looks at" the encoded input to decide what information is relevant.

\textbf{Computation:}

\begin{equation}
\text{attended}_{ij} = \sum_k \alpha_{ijk} \cdot V_{jk}
\end{equation}

Where:

\begin{equation}
\alpha_{ijk} = \frac{\exp\left(\frac{Q_{id} K_{jd}^T}{\sqrt{D}}\right)}{\sum_m \exp\left(\frac{Q_{id} K_{md}^T}{\sqrt{D}}\right)}
\end{equation}

\textbf{Output:} $\text{Attended Slots} \in \mathbb{R}^{16 \times 21 \times 64}$

\subsection{Self-Attention Between Slots}

\textbf{Explanation:}

Now the 21 slots communicate with each other to specialize.

\textbf{Computation (same multi-head attention):}

\begin{equation}
\text{refined}_{ij} = \text{MultiHeadAttention}(\text{attended}_i, \text{attended}_i, \text{attended}_i)
\end{equation}

\textbf{Output:} $\text{Refined Slots} \in \mathbb{R}^{16 \times 21 \times 64}$

\subsection{MLP H0: Decode Slot 1}

\textbf{Input:} refined\_slots[:, 0, :] (first slot) $\in \mathbb{R}^{16 \times 64}$

\textbf{MLP:}

\begin{equation}
h = \text{ReLU}(W_1 \cdot \text{slot0} + b_1)
\end{equation}

\begin{equation}
\text{h0\_flat} = W_2 \cdot h + b_2
\end{equation}

\textbf{Reshape to matrix:}

\begin{equation}
H_0 = \text{reshape}(\text{h0\_flat}, (16, 3, 3))
\end{equation}

\textbf{Output:} $H_0 \in \mathbb{R}^{16 \times 3 \times 3}$

\subsection{MLP Drops: Decode Slots 2-21}

\textbf{Input:} refined\_slots[:, n, :] for $n = 1, \ldots, 20$ each in $\mathbb{R}^{16 \times 64}$

\textbf{MLP for each slot:}

\begin{equation}
h_n = \text{ReLU}(W_1 \cdot \text{slot}_n + b_1)
\end{equation}

\begin{equation}
\text{output}_n = W_2 \cdot h_n + b_2 \in \mathbb{R}^{16 \times 10}
\end{equation}

\textbf{Split output:}

\begin{equation}
H_n = \text{reshape}(\text{output}_n[:, :9], (16, 3, 3))
\end{equation}

\begin{equation}
w_n = \sigma(\text{output}_n[:, 9]) = \frac{1}{1 + e^{-\text{output}_n[:, 9]}} \in [0, 1]
\end{equation}

\section{Loss Functions and Training}

\subsection{Data Loss}

\textbf{Formula:}

\begin{equation}
\mathcal{L}_{\text{data}} = \frac{1}{B}\sum_{b=1}^{B} (f_{\text{pred}}^b - f_{\text{true}}^b)^2
\end{equation}

Where:
\begin{itemize}
  \item $f_{\text{pred}}^b$: predicted frequency for sample $b$
  \item $f_{\text{true}}^b$: ground truth frequency for sample $b$
  \item $B = 16$: batch size
\end{itemize}

\textbf{Example:}
\begin{lstlisting}
pred_freq = [2.3, 2.1, 1.9, ...]  # 16 predicted values
true_freq = [2.2, 2.0, 2.1, ...]  # 16 ground truth values
data_loss = mean((pred_freq - true_freq)^2)
\end{lstlisting}

\subsection{Frequency Consistency Loss (LoT1)}

\textbf{Formula:}

\begin{equation}
\mathcal{L}_{\text{freq}} = \frac{1}{B}\sum_{b=1}^{B} (f_{\text{pred}}^b - f_{\text{analytical}}^b)^2
\end{equation}

Where:
\begin{itemize}
  \item $f_{\text{analytical}}^b$: analytically computed frequency (from FEM)
  \item Enforces: predicted frequency matches physics-based calculation
\end{itemize}

\subsection{Monotonicity Loss (LoT2)}

\textbf{Formula:}

\begin{equation}
\mathcal{L}_{\text{mono}} = \frac{1}{B}\sum_{b=1}^{B} \sum_{n=1}^{19} \max(0, w_n^b - w_{n+1}^b)
\end{equation}

Where:
\begin{itemize}
  \item $w_n^b$: activation gate for sample $b$, drop $n$
  \item $\max(0, \cdot)$: penalty only if gate is decreasing (which is bad)
\end{itemize}

\textbf{Purpose:} Enforce $w_1 \leq w_2 \leq \cdots \leq w_{20}$ (monotonic increase)

\subsection{Positive Stiffness Loss (LoT2)}

\textbf{Formula:}

\begin{equation}
\mathcal{L}_{\text{pos}} = \frac{1}{B}\sum_{b=1}^{B} \sum_{i=1}^{3} \max(0, -\lambda_i(H_{\text{evolved}}^b))
\end{equation}

Where:
\begin{itemize}
  \item $\lambda_i(H)$: i-th eigenvalue of matrix $H$
  \item $\max(0, -\lambda)$: penalty if eigenvalue is negative
\end{itemize}

\textbf{Purpose:} Ensure all eigenvalues are non-negative (positive definite matrix)

\subsection{Composite Loss}

\textbf{Formula:}

\begin{equation}
\mathcal{L}_{\text{total}} = \alpha \mathcal{L}_{\text{data}} + \beta \mathcal{L}_{\text{freq}} + \gamma \mathcal{L}_{\text{mono}} + \delta \mathcal{L}_{\text{pos}}
\end{equation}

Where weights are:
\begin{itemize}
  \item $\alpha = 1.0$: Emphasize data fitting
  \item $\beta = 0.5$: Enforce physics consistency
  \item $\gamma = 0.3$: Enforce smooth degradation
  \item $\delta = 0.2$: Prevent stiffness collapse
\end{itemize}

\subsection{Backpropagation}

\textbf{Chain of derivatives:}

\begin{equation}
\frac{\partial \mathcal{L}}{\partial \theta} = \frac{\partial \mathcal{L}}{\partial \text{freq\_pred}} \cdot \frac{\partial \text{freq\_pred}}{\partial \text{eq\_token}} \cdot \cdots \cdot \frac{\partial H_0}{\partial \theta}
\end{equation}

\textbf{Parameter update (Adam optimizer):}

\begin{equation}
\theta_{\text{new}} = \theta_{\text{old}} - \text{lr} \cdot m_t / \sqrt{v_t + \epsilon}
\end{equation}

Where:
\begin{itemize}
  \item $m_t$: first moment estimate (momentum)
  \item $v_t$: second moment estimate (variance)
  \item $\text{lr} = 0.001$: learning rate
  \item $\epsilon = 10^{-8}$: small constant for stability
\end{itemize}

\section{Complete Data Flow Example}

\textbf{One complete forward pass:}

\begin{enumerate}
  \item \textbf{Input}: 16 load sequences, 16 soil profiles
  \item \textbf{Step 1}: Embed loads $(16, 50, 1) \rightarrow (16, 50, 64)$
  \item \textbf{Step 2}: Slot-Attention on soil params $(16, 8) \rightarrow (H_0, H_d, w_g)$
  \item \textbf{Step 3}: Evolve stiffness $\rightarrow H_{\text{evolved}} (16, 3, 3)$
  \item \textbf{Step 4}: Concat loads + stiffness $\rightarrow (16, 51, 64)$
  \item \textbf{Step 5}: Encoder processes all tokens $\rightarrow (16, 51, 64)$
  \item \textbf{Step 6}: Decoder cross-attends $\rightarrow (16, 51, 64)$
  \item \textbf{Step 7}: Extract freq token, predict $\rightarrow (16,)$
  \item \textbf{Output}: 16 predicted frequencies
\end{enumerate}

\textbf{Then compute losses, backprop, and update all weights.}

\section{Summary Table}

\begin{table}[H]
\centering
\begin{tabular}{|l|l|l|}
\hline
\textbf{Component} & \textbf{Input Shape} & \textbf{Output Shape} \\
\hline
Load Embedding & $(16, 50, 1)$ & $(16, 50, 64)$ \\
\hline
Slot-Attention & $(16, 8)$ & $H_0: (16,3,3)$, $H_d: (16,20,3,3)$, $w_g: (16,20)$ \\
\hline
Stiffness Evolution & $H_0, H_d, w_g$ & $(16, 3, 3)$ \\
\hline
Concatenation & Loads + stiffness & $(16, 51, 64)$ \\
\hline
Encoder & $(16, 51, 64)$ & $(16, 51, 64)$ \\
\hline
Decoder & $(16, 51, 64)$ & $(16, 51, 64)$ \\
\hline
Freq Head & $(16, 64)$ & $(16,)$ \\
\hline
\end{tabular}
\end{table}

\end{document}
